% !TEX root = ../main.tex
\section{Introduction}
\label{sec:intro}

\IEEEPARstart{V}{ideo} Diffusion Transformers (VDiTs) enable high-quality generation but repeatedly execute expensive Spatial-Temporal Diffusion Transformer (STDiT) blocks during denoising. Open-Sora v1.2, for example, evaluates 28 spatial--temporal block pairs over 30 sampling steps~\cite{zheng2024opensora}. Timestep caching skips redundant STDiT evaluations~\cite{liu2025teacache}, whereas post-training quantization reduces the cost of each recomputed step~\cite{zhao2025viditq}. Recent joint methods further optimize cache schedules, correction, or precision selection~\cite{liu2025cachequant,wu2025quantcache}, while dedicated VDiT acceleration such as ViDA exploits differential approximation and adaptive dataflow~\cite{ding2025vida}. These optimizations attack different dimensions of cost, but they become coupled in a cache--quantized deployment: the cache policy determines which denoising steps reach the accelerator, while quantization and dataflow determine both the numerical trajectory and the execution cost of those recomputed steps. Directly composing them therefore leaves two cross-layer mismatches.

First, a \textbf{decision-domain mismatch} arises when a full-precision (FP) cache indicator ranks timestep changes that are later altered by quantization. Quantization is nonlinear and piecewise constant: an FP difference can disappear inside one quantization bin, while a smaller change can cross a bin boundary. A policy calibrated only on FP variation can therefore mis-rank the benefit of recomputation once the actual execution path is quantized. This mismatch matters particularly for a fixed W8A8/INT8 accelerator, where the cache controller cannot rely on a separate FP recomputation path to recover the original indicator semantics.

Second, a \textbf{matrix-axis mismatch} arises because Open-Sora's factorized spatial/temporal attention~\cite{zheng2024opensora} places the varying sequence dimension on different GEMM axes. In $QK^{\mathsf T}$ it determines the output width; in $PV$ it becomes the reduction length. Our 64-, 256-, and 512-frame workloads have temporal lengths $T=19$, 76, and 151, respectively. The $T=19$ workload is a representative short-axis case: mapping one head to all 36 lanes fragments the $QK^{\mathsf T}$ output, while row-wise padding wastes the same array in $PV$. A single static lane/reduction policy therefore cannot efficiently cover both GEMMs, and duplicating a temporal-specialized array would sacrifice the area advantage expected from an accelerator.

Fig.~\ref{fig:overview} summarizes \textbf{Q-DARE}, which addresses these mismatches at the two layers where they arise. \textbf{Quantization-aligned caching (QAC)} evaluates cache relevance in the same numerical domain as recomputation and calibrates the indicator against the quantized trajectory, aligning \emph{which} steps are recomputed. The \textbf{dimension-adaptive reconfigurable engine (DARE)} improves \emph{how} those steps are executed by reusing one processing-element array (PEA): temporal $QK^{\mathsf T}$ changes output-lane ownership, whereas temporal $PV$ packs boundary-tagged reductions and replays V without an explicit transpose. These mechanisms are complementary rather than alternative: QAC reduces unnecessary STDiT executions, while DARE increases utilization of the remaining ones without adding another precision path or another PEA.

The resulting co-design preserves generation quality while improving both GPU-relative and accelerator-relative performance. Q-DARE reaches 76.61 VBench versus 76.77 for FP16. For the representative $T=19$ case, DARE raises temporal $QK^{\mathsf T}$ lane utilization from 52.8\% to 70.4\% and temporal $PV$ product utilization from 52.8\% to 91.2\%. Across 64--512-frame workloads, Q-DARE achieves \QDAREStepGeoA$\times$ geomean STDiT step-throughput speedup and \QDAREEnergyGeoA$\times$ energy-efficiency improvement over A100, while reaching \QDAREViDAThroughputGeo$\times$ geomean video throughput and \QDAREViDAreaGeo$\times$ geomean area efficiency over ViDA.

Our contributions are:
\begin{itemize}
  \item We propose QAC, a quantization-aligned cache indicator that makes reuse decisions in the same numerical domain as fixed W8A8/INT8 recomputation.
  \item We design DARE, a shared datapath that adapts the output and reduction axes of factorized attention without duplicating the PEA.
  \item We implement Q-DARE in RTL and demonstrate near-FP16 quality with large efficiency gains.
\end{itemize}
